%%%%%%%%%%%%%%%%%%%%%%preface.tex%%%%%%%%%%%%%%%%%%%%%%%%%%%%%%%%%%%%%%%%%
% sample preface
%
% Use this file as a template for your own input.
%
%%%%%%%%%%%%%%%%%%%%%%%% Springer %%%%%%%%%%%%%%%%%%%%%%%%%%

%\preface

\chapter*{Przedmowa}
\thispagestyle{empty}
% \addcontentsline{toc}{chapter}{Przedmowa}

%Todo co w jakim rozdziale
Niniejszy tom prezentuje prace naukowe obejmujące szeroki zakres zagadnień związanych ze współczesną informatyką. Zgromadzone teksty podejmują tematy istotne dla rozwoju tej dyscypliny oraz obszarów z nią powiązanych, ukazując jej interdyscyplinarny charakter i rosnące znaczenie we współczesnej nauce i~technologii.

W publikacji znalazły się opracowania dotyczące m.in. telekomunikacji, której dynamiczny rozwój umożliwia tworzenie coraz bardziej zaawansowanych systemów komunikacji; matematyki, stanowiącej podstawę modelowania, analizy i~projektowania algorytmów; informatyki kwantowej, otwierającej nowe kierunki badań w obliczeniach i kryptografii; oraz sztucznej inteligencji, będącej jednym z~kluczowych motorów innowacji technologicznych ostatnich lat. Zestawienie tych obszarów w jednym tomie podkreśla, jak szeroko i wielowymiarowo można ujmować współczesne badania informatyczne.

Celem publikacji jest przedstawienie aktualnych trendów, podejść badawczych i wyników, które przyczyniają się do lepszego zrozumienia zarówno teoretycznych, jak i praktycznych aspektów informatyki. Mamy nadzieję, że różnorodność omawianych zagadnień stanie się dla Czytelnika źródłem inspiracji oraz zachętą do dalszych poszukiwań i pogłębiania wiedzy.

Tom kierowany jest do badaczy, praktyków oraz studentów zainteresowanych rozwojem informatyki i jej zastosowaniami, oferując przekrojowe spojrzenie na wybrane nurty badań prowadzonych w tej dziedzinie.

 

\vspace{\baselineskip}
\begin{flushright}\noindent
Wroc\l aw, Grudzień 2025\hfill {\it Radosław Idzikowski i Mateusz A. Kucharski}\\
\end{flushright}


